% ID: SEL-D0031-OB-E03
% ESTADO: cotejada. Segunda pasada visual completa y borrador compilado revisados.
% Fuente: Selectividad/matematicas_aplicadas/titular_junio_6 EXAMEN MAT APLIC C S.pdf
% SHA-256: bebc339cfd8f69c7c4fae160fd165eaef5f6bc93cfd6d15cb9ce289b7029d167
% Página: 2; opción B; ejercicio original 3.

En una encuesta sobre la nacionalidad de los veraneantes en un municipio de la costa andaluza, se ha observado que el 40\% de los encuestados son españoles y el 60\% extranjeros, que el 30\% de los españoles y el 80\% de los extranjeros residen en un hotel y el resto en otro tipo de residencia.

Se elige al azar un veraneante del municipio.
\begin{enumerate}[label=\alph*)]
\item \textbf{(1 punto)} \textquestiondown{}Cuál es la probabilidad de que no resida en un hotel?
\item \textbf{(1 punto)} Si no reside en un hotel, \textquestiondown{}cuál es la probabilidad de que sea español?
\item \textbf{(0.5 puntos)} \textquestiondown{}Son independientes los sucesos ``ser extranjero'' y ``residir en un hotel''?
\end{enumerate}
